\documentclass[11pt]{article}
\usepackage[a4paper,margin=26mm]{geometry}
\usepackage[T1]{fontenc}
\usepackage{lmodern}
\usepackage{amsmath,amssymb}
\usepackage[hidelinks]{hyperref}

\title{Maximum Tail Length One for $x^2+1$\\[4pt]
  \large Disproof of Conjecture 00000007756}
\author{Feiyu (Independent researcher)}
\date{26 November 2026}

\begin{document}
\maketitle

\begin{abstract}
For the functional graph of $x\mapsto x^2+1$ over $\mathbb F_p$, the
maximum tail length is one exactly at $p=5$. Its frequency over the
primes tends to zero. If occurrence means that the graph contains a
vertex of tail length one, the frequency tends to one. Both limits
differ from the conjectured $3-4\log 2$.
\end{abstract}

\section{The tail-frequency claim}

The \href{https://github.com/The-Last-Math-Competition/The-Last-Math-Competition/blob/main/conjectures/00000007756.md}{bilingual
statement} fixes $f(x)=x^2+1$ and varies the prime $p$, with one
functional graph per prime. Its first assertion gives the limiting
frequency of shapes of tail length one as $3-4\log 2$. The second
assertion concerns the distribution of the longest cycle scaled by $p$.

Write $f_p(x)=x^2+1$ on $\mathbb F_p$. The graph has the edge
$x\to f_p(x)$ at every vertex. The tail length $t_p(x)$ is the least
$t\ge0$ for which $f_p^t(x)$ is on a directed cycle \cite[Section~1]{vs}.
Put $T_p=\max_{x\in\mathbb F_p} t_p(x)$. A graph has maximum tail length
one if and only if every $f_p(x)$ is cyclic and some $x$ is noncyclic.
The asserted limit for this graph property is
\[
  \lim_{X\to\infty}
  \frac{\#\{p\le X : p\text{ prime},\ T_p=1\}}{\pi(X)}
  =3-4\log 2 ,
\]
where $\pi(X)$ counts primes at most $X$ and $\log$ is the natural
logarithm. Section~\ref{sec:frequency} also treats the graph property
$\exists x\in\mathbb F_p$ with $t_p(x)=1$.

\section{An obstruction for every prime above five}

Fix a prime $p>5$, write $I=f_p(\mathbb F_p)$, and suppose all points
of $I$ are cyclic. The map on cyclic points is a permutation, so it
is injective on $I$. Since $f_p(z)=f_p(-z)$, any $z\in I\cap(-I)$
must satisfy $z=-z$, hence $z=0$.

The residues represented by $0,1,\ldots,(p-1)/2$ have distinct images.
Indeed, if $f_p(a)=f_p(b)$, then $a=b$ or $a=-b$ in $\mathbb F_p$.
The first case forces equality of the representatives. In the second,
$a+b$ is divisible by $p$, and $0\le a+b\le p-1$, so both are zero.
Thus $|I|\ge (p+1)/2$. The sets $I$ and $-I$ intersect since their
cardinalities sum to more than $p$. Therefore $0\in I$, and there
exists $i\in\mathbb F_p$ satisfying $i^2=-1$.

Now put $u=1/2$ and $v=3i/2$. Their images are
\[
  f_p(u)=\frac54,\qquad f_p(v)=-\frac54 .
\]
These are nonzero because $p>5$, contradicting $I\cap(-I)\subseteq\{0\}$.
Consequently there is a vertex $x$ whose successor $f_p(x)$ is
noncyclic. This gives $t_p(x)\ge2$.

\newpage

\section{Exact classification and limiting frequencies}
\label{sec:frequency}

The prime fields of size at most five give the following graphs.
In the table the tail lengths are in the order of the residues
$0,1,\ldots,p-1$.
\[
\begin{array}{c|l|c}
p & \text{Cyclic part} & (t_p(0),\ldots,t_p(p-1)) \\ \hline
2 & 0\to1\to0 & (0,0) \\
3 & 2\to2 & (2,1,0) \\
5 & 0\to1\to2\to0 & (0,0,0,1,1)
\end{array}
\]
For $p=3$, $0\to1\to2$. For $p=5$, the two remaining edges are
$3\to0$ and $4\to2$. Together with the uniform obstruction, this
shows that $T_p=1$ exactly at $p=5$. For every $X\ge5$ the frequency
is $1/\pi(X)$, which tends to $0$ by the infinitude of primes.

For the graph property of containing a tail of length one, take any
prime $p>2$. The vertices $1$ and $-1$ are distinct and have the same
successor, so at least one is noncyclic. A forward orbit of a
noncyclic vertex in a finite graph eventually enters a cycle. The
vertex immediately before its first entry has tail length one.
At $p=2$ all vertices are cyclic. Thus for every $X\ge2$ this
frequency is $1-1/\pi(X)$, which tends to $1$.

On $1<t<2$ one has $1/2<1/t<(3-t)/2$. Integrating gives
$1/2<\log 2<3/4$, so $0<3-4\log 2<1$. The tail-frequency assertion
is false under either of the graph properties above, which disproves
the conjunction with the cycle-length assertion.

\section{Lean formalization}

The project uses Lean and Mathlib v4.33.1. The map
\texttt{quadraticMap} acts on \texttt{ZMod p}, and
\texttt{Function.periodicPts} gives the set of cyclic vertices.
\texttt{MaxTailOne} expresses $T_p=1$ through cycle membership after
one step and the existence of a noncyclic vertex.
\texttt{ContainsTailOne} checks for a vertex at distance one from a cycle.

\begin{itemize}
\item \texttt{exists\_successor\_off\_cycle} proves the obstruction for
  every prime $p>5$ using an injective set of representatives,
  finite-set cardinalities and the opposite pair of nonzero images.
\item \texttt{maxTailOne\_iff\_eq\_five} combines this with the three
  small cases. For the finite checks, a vertex is cyclic exactly
  when it returns after some positive number of steps at most $p$.
\item The two limit theorems use Mathlib's prime-counting function.\newline
  \texttt{tailOneFrequency\_tendsto\_zero}\newline
  \texttt{containsTailOneFrequency\_tendsto\_one}
\item \texttt{conjecture\_00000007756\_false} negates the conjunction
  of the primary frequency claim with any second proposition.
  \texttt{not\_containsTailFrequencyClaim} refutes the alternative
  occurrence reading. The axiom audits use only \texttt{propext},
  \texttt{Classical.choice} and \texttt{Quot.sound}.
\end{itemize}

\begin{thebibliography}{9}
\bibitem{vs}
T. Vasiga and J. Shallit,
\href{https://cs.uwaterloo.ca/~shallit/Papers/vas.pdf}
     {On the iteration of certain quadratic maps over $\mathrm{GF}(p)$},
\emph{Discrete Mathematics} \textbf{277} (2004), 219--240.
\end{thebibliography}

\end{document}
